% notes/artefacts/model-governance-controls/derivations.tex
% Legacy blocks kept verbatim when they were migrated (notes/migration.md), for the pathway
% model-governance-controls and the entries that took blocks from the same legacy file. The compendium
% entries state each claim with its locator; the original wording lives here.
% Treat everything here as unchecked own working.
% Compile inside notes/artefacts/model-governance-controls/ with  latexmk -pdf derivations.tex
\PassOptionsToPackage{table}{xcolor}
\documentclass[a4paper,oneside,openany]{book}
% Your own packages, macros and environments. The learning guide's entry environments,
% \citep and \qfsource are in house.tex, which compendium.tex loads after this file.

% -----------------------------
% Document Class & Encoding
% -----------------------------
\usepackage[utf8]{inputenc}  % Input encoding
\usepackage[T1]{fontenc}     % Font encoding

% -----------------------------
% Page Layout & Geometry
% -----------------------------
\usepackage[margin=2cm,foot=1cm]{geometry}

% -----------------------------
% Header & Footer Management
% -----------------------------
\usepackage{fancyhdr}
\pagestyle{fancy}

% -----------------------------
% Text Formatting & Spacing
% -----------------------------
\usepackage{ragged2e}
\usepackage[nodisplayskipstretch]{setspace}
\setstretch{1}  % Line spacing

% -----------------------------
% Listings for Code
% -----------------------------
\usepackage{listings}
\usepackage{xcolor}

% -----------------------------
% Tables, Figures & Graphics
% -----------------------------
\usepackage{tabularx}
\usepackage{multirow}
\usepackage{pdflscape}
\usepackage{graphicx}

% -----------------------------
% Algorithms
% -----------------------------
\usepackage{algorithm}       % For floating algorithm environment
\usepackage{algpseudocode}   % For pseudo-code inside algorithms

% -----------------------------
% Math Packages
% -----------------------------
\usepackage{amsmath,amssymb,amsthm,mathtools,bm}
\usepackage{mathrsfs}        % Script fonts
\usepackage{tikz-cd}         % Commutative diagrams
\usepackage{rotating}        % Rotated objects
\usepackage{float}           % Float placement
\usetikzlibrary{calc, arrows, positioning, arrows.meta}

% -----------------------------
% TikZ and its Libraries
% -----------------------------
\usepackage{tikz}
\usetikzlibrary{positioning} % Added to support "below=of" syntax
\usepackage{forest}          % For advanced tree diagrams
\usepackage{tikz-qtree}      % For vertical tree diagrams

% -----------------------------
% Miscellaneous Packages
% -----------------------------
\usepackage{xparse}          % For defining robust commands
\usepackage{enumitem}        % For customizing list layouts
\usepackage{makecell}        % For enhanced table cells
\usepackage{comment}		 % For commenting out sections

%%%%%%%%%%%%%%%%%%%%%%%%%%%%%%%%%%%%%%%%%%%%%%%%%%%%%%%%%%%%%%%%%%%%%%%%%%%%%%%
% Custom Math Commands and Shortcuts
%%%%%%%%%%%%%%%%%%%%%%%%%%%%%%%%%%%%%%%%%%%%%%%%%%%%%%%%%%%%%%%%%%%%%%%%%%%%%%%

% Matrix and Vector Formatting
\newcommand{\matr}[1]{\mathbf{#1}}      % For undergraduate algebra (alternative: \bm{#1})
\newcommand{\vect}[1]{\bm{#1}}

% Standard Number Sets
\newcommand{\N}{\mathbb{N}}
\newcommand{\R}{\mathbb{R}}
\newcommand{\Q}{\mathbb{Q}}
\newcommand{\C}{\mathbb{C}}
\newcommand{\Z}{\mathbb{Z}}

% Bold variables (customize as needed)
\newcommand{\x}{\mathbf{x}}
\newcommand{\F}{\mathbf{F}}
\newcommand{\f}{\mathbf{f}}
\newcommand{\y}{\mathbf{y}}
\renewcommand{\a}{\mathbf{a}}
\renewcommand{\b}{\mathbf{b}}
\renewcommand{\c}{\mathbf{c}}
\newcommand{\h}{\mathbf{h}}
\newcommand{\g}{\mathbf{g}}
\newcommand{\z}{\mathbf{z}}
\newcommand{\ze}{\mathbf{0}}

% Norms, Absolute Values, and Bracketing
\newcommand{\norm}[1]{\left\lVert#1\right\rVert}
\newcommand{\abs}[1]{\left\lvert#1\right\rvert}
\newcommand{\brk}[1]{\left[#1\right]}
\newcommand{\brc}[1]{\left\{#1\right\}}
\newcommand{\paren}[1]{\left(#1\right)}
\newcommand{\normop}[1]{\left\lVert#1\right\rVert_\text{op}}

% Miscellaneous Math Shortcuts
\newcommand{\LL}{\mathcal{L}}
\newcommand{\uni}{\overset{\text{uni}}{\to}}
\DeclareMathOperator{\diam}{diam}
\newcommand{\Prr}[1]{\text{Pr}\left(#1\right)}
\newcommand{\ceil}[1]{\lceil #1 \rceil}
\newcommand{\floor}[1]{\lfloor #1 \rfloor}
\newcommand\disteq{\stackrel{\textnormal{dist}}{=}}

% Evaluation command with options
\NewDocumentCommand{\evalat}{sO{\big}mm}{%
  \IfBooleanTF{#1}
    {\mleft. #3 \mright|_{#4}}
    {#3#2|_{#4}}%
}

%%%%%%%%%%%%%%%%%%%%%%%%%%%%%%%%%%%%%%%%%%%%%%%%%%%%%%%%%%%%%%%%%%%%%%%%%%%%%%%
% Color Definitions
%%%%%%%%%%%%%%%%%%%%%%%%%%%%%%%%%%%%%%%%%%%%%%%%%%%%%%%%%%%%%%%%%%%%%%%%%%%%%%%

\definecolor{dartmouthgreen}{rgb}{0.05, 0.5, 0.06}
\definecolor{egyptianblue}{rgb}{0.06, 0.2, 0.65}
\definecolor{dukeblue}{rgb}{0.0, 0.0, 0.61}
\definecolor{jazzberryjam}{rgb}{0.65, 0.04, 0.37}
\definecolor{magenta}{HTML}{EC008C}
\definecolor{darkmagenta}{rgb}{0.55, 0.0, 0.55}
\definecolor{deeppink}{rgb}{1.0, 0.08, 0.58}
\definecolor{codegreen}{rgb}{0,0.6,0}
\definecolor{codegray}{rgb}{0.5,0.5,0.5}
\definecolor{codepurple}{rgb}{0.58,0,0.82}
\definecolor{backcolour}{rgb}{0.95,0.95,0.92}

%%%%%%%%%%%%%%%%%%%%%%%%%%%%%%%%%%%%%%%%%%%%%%%%%%%%%%%%%%%%%%%%%%%%%%%%%%%%%%%
% Custom Shortcuts and Emphasis Commands
%%%%%%%%%%%%%%%%%%%%%%%%%%%%%%%%%%%%%%%%%%%%%%%%%%%%%%%%%%%%%%%%%%%%%%%%%%%%%%%

\newcommand{\isomorp}{\xrightarrow{\sim}}
\newcommand{\hlt}[1]{\textit{{\color{blue}#1}}}
\newcommand{\simpt}[1]{\textit{{\color{deeppink}#1}}}
\newcommand{\impt}[1]{\textit{{\color{red}#1}}}
\newcommand{\gcds}[1]{\textnormal{gcd}#1}
\newcommand{\mins}[1]{\textnormal{min}#1}
\newcommand{\maxs}[1]{\textnormal{max}#1}
\newcommand{\lcms}[1]{\textnormal{lcm}#1}
\newcommand{\degs}[1]{\textnormal{deg}#1}
\newcommand{\exps}[1]{\textnormal{exp}#1}
\newcommand{\tors}[1]{\textnormal{Tor}#1}
\newcommand{\homs}[1]{\textnormal{Hom}#1}
\newcommand{\anns}[1]{\textnormal{Ann}#1}
\newcommand{\defeq}{\vcentcolon=}
\DeclareMathOperator{\sign}{sign}

%%%%%%%%%%%%%%%%%%%%%%%%%%%%%%%%%%%%%%%%%%%%%%%%%%%%%%%%%%%%%%%%%%%%%%%%%%%%%%%
% Theorem and Definition Environments
%%%%%%%%%%%%%%%%%%%%%%%%%%%%%%%%%%%%%%%%%%%%%%%%%%%%%%%%%%%%%%%%%%%%%%%%%%%%%%%
\theoremstyle{plain}
\newtheorem{theorem}{Theorem}[section]
\newtheorem{corollary}[theorem]{Corollary}
\newtheorem{lemma}[theorem]{Lemma}
\newtheorem{proposition}[theorem]{Proposition}
\newtheorem{qnbank}[theorem]{Question Bank}
\let\oldqnbank\qnbank
\renewcommand{\qnbank}{\oldqnbank\normalfont}
\newtheorem{process}[theorem]{Process}
\let\oldprocess\process
\renewcommand{\process}{\oldprocess\normalfont}
\newtheorem{method}[theorem]{Method}
\let\oldmethod\method
\renewcommand{\method}{\oldmethod\normalfont}
\newtheorem{definition}[theorem]{Definition}
\let\olddefinition\definition
\renewcommand{\definition}{\olddefinition\normalfont}
\newtheorem{example}[theorem]{Example}
\let\oldexample\example
\renewcommand{\example}{\oldexample\normalfont}
\newtheorem{remark}[theorem]{Remark}
\let\oldremark\remark
\renewcommand{\remark}{\oldremark\normalfont}
\newtheorem{notation}[theorem]{Notation}
\let\oldnotation\notation
\renewcommand{\notation}{\oldnotation\normalfont}
\newtheorem{principle}[theorem]{Principle}
\let\oldprinciple\principle
\renewcommand{\principle}{\oldprinciple\normalfont}
\newtheorem{mechanism}[theorem]{Mechanism}
\let\oldmechanism\mechanism
\renewcommand{\mechanism}{\oldmechanism\normalfont}
\newtheorem{taxonomy}[theorem]{Taxonomy}
\let\oldtaxonomy\taxonomy
\renewcommand{\taxonomy}{\oldtaxonomy\normalfont}
\newtheorem{model}[theorem]{Model}
\let\oldmodel\model
\renewcommand{\model}{\oldmodel\normalfont}
\newtheorem{assumption}[theorem]{Assumption}
\let\oldassumption\assumption
\renewcommand{\assumption}{\oldassumption\normalfont}
\newtheorem{casestudy}[theorem]{Case Study}
\let\oldcasestudy\casestudy
\renewcommand{\casestudy}{\oldcasestudy\normalfont}
\newtheorem{empiricalfinding}[theorem]{Empirical Finding}
\let\oldempiricalfinding\empiricalfinding
\renewcommand{\empiricalfinding}{\oldempiricalfinding\normalfont}
\newtheorem{practitionernote}[theorem]{Practitioner Note}
\let\oldpractitionernote\practitionernote
\renewcommand{\practitionernote}{\oldpractitionernote\normalfont}
\newtheorem{caveat}[theorem]{Caveat}
\let\oldcaveat\caveat
\renewcommand{\caveat}{\oldcaveat\normalfont}
\newtheorem{marketrule}[theorem]{Rule}
\let\oldmarketrule\marketrule
\renewcommand{\marketrule}{\oldmarketrule\normalfont}
\newtheorem{takeaway}[theorem]{Takeaway}
\let\oldtakeaway\takeaway
\renewcommand{\takeaway}{\oldtakeaway\normalfont}

%%%%%%%%%%%%%%%%%%%%%%%%%%%%%%%%%%%%%%%%%%%%%%%%%%%%%%%%%%%%%%%%%%%%%%%%%%%%%%%
% Breakable Algorithm Environment (Algorithms spanning multiple pages)
%%%%%%%%%%%%%%%%%%%%%%%%%%%%%%%%%%%%%%%%%%%%%%%%%%%%%%%%%%%%%%%%%%%%%%%%%%%%%%%

\makeatletter
\newenvironment{breakablealgorithm}
 {%
  \begin{center}
    \refstepcounter{algorithm} % Increment algorithm counter
    \hrule height.8pt depth0pt \kern2pt
    \renewcommand{\caption}[2][\relax]{%
      {\raggedright\textbf{Algorithm \thealgorithm:} ##2\par}%
      \ifx\relax##1\relax
        \addcontentsline{loa}{algorithm}{\protect\numberline{\thealgorithm}##2}%
      \else
        \addcontentsline{loa}{algorithm}{\protect\numberline{\thealgorithm}##1}%
      \fi
      \kern2pt\hrule\kern2pt
    }%
 }{%
    \kern2pt\hrule
  \end{center}
 }
\makeatother

%%%%%%%%%%%%%%%%%%%%%%%%%%%%%%%%%%%%%%%%%%%%%%%%%%%%%%%%%%%%%%%%%%%%%%%%%%%%%%%
% Code Listings Style
%%%%%%%%%%%%%%%%%%%%%%%%%%%%%%%%%%%%%%%%%%%%%%%%%%%%%%%%%%%%%%%%%%%%%%%%%%%%%%%

\lstdefinestyle{codestyle}{
  commentstyle=\color{codegreen},
  keywordstyle=\color{magenta},
  numberstyle=\tiny\color{codegray},
  stringstyle=\color{codepurple},
  basicstyle=\ttfamily\normalsize,
  breakatwhitespace=false,
  breaklines=true,
  captionpos=b,
  keepspaces=true,
  numbers=left,
  numbersep=10pt,
  frame=single,
  showspaces=false,
  showstringspaces=false,
  showtabs=false,
  tabsize=2
}
\lstset{style=codestyle}

% -----------------------------
% Hyperlinks (loaded last: loaded before float and algorithm it duplicated figure anchors)
% -----------------------------
\usepackage{hyperref}
\hypersetup{
  colorlinks,
  linkcolor={blue!80!black},
  citecolor={green!80!black},
  urlcolor={red!80!black}
}

\usepackage[T1]{fontenc}
\usepackage{amsmath,amssymb,booktabs,graphicx,hyperref}
\hypersetup{colorlinks}
\setlength{\emergencystretch}{3em}   % lets a line holding a long key break
\newcounter{qfentry}[chapter]
\renewcommand{\theqfentry}{\thechapter.\arabic{qfentry}}
% \qftype{type}{Name} declares the environment qf<type>; its title wraps like a paragraph.
\newcommand{\qftype}[2]{\newenvironment{qf#1}[1][]{\par\addvspace{\medskipamount}%
  \refstepcounter{qfentry}{\raggedright\noindent\bfseries #2 \theqfentry
  \if\relax\detokenize{##1}\relax\else\quad ##1\fi\par}\nopagebreak}%
  {\par\addvspace{\medskipamount}}}
\qftype{definition}{Definition}\qftype{convention}{Convention}\qftype{result}{Result}
\qftype{model}{Model}\qftype{method}{Method}\qftype{mechanism}{Mechanism}
\qftype{evidence}{Evidence}\qftype{data}{Data}\qftype{rule}{Rule}
\qftype{pitfall}{Pitfall}\qftype{example}{Example}
\newcommand{\qffield}[1]{\par\textbf{#1.}\enspace}   % starts one of an entry's fields
% \citep[locator]{key} links the key to its line in the sources if it has one,
% and otherwise to its entry in the pinned catalogue PDF.
\newcommand{\qfcatalogue}{../guide/institutional_quantitative_finance_reading_list.pdf}
\NewDocumentCommand{\citep}{o o m}{[\IfValueT{#2}{#1 }%
  \ifcsname r@src:#3\endcsname\hyperref[src:#3]{\nolinkurl{#3}}%
  \else\href{\qfcatalogue\##3}{\nolinkurl{#3}}\fi
  \IfValueTF{#2}{, #2}{\IfValueT{#1}{, #1}}]}
% \qfsource{key}{reference} records one source outside the catalogue, in front/sources.tex.
\newcommand{\qfsource}[2]{\par\noindent\phantomsection\label{src:#1}\nolinkurl{#1}\quad #2\par}
% \qfsections{key, ...} inputs <key>/section.tex for each listed key that has one, in order.
\ExplSyntaxOn
\NewDocumentCommand{\qfsections}{m}
  {\clist_map_inline:nn {#1} {\InputIfFileExists{##1/section}{}{}}}
\ExplSyntaxOff

\pagestyle{plain}
\begin{document}
\chapter*{Model governance and trading controls: legacy derivations}

\section*{Production considerations (slice 17, 2026-10-15)}
% Kept verbatim when this file was migrated on 2026-10-15 (notes/migration.md, slice 17): the wrapper
% of the folder 01_execution_and_algorithmic_trading (00_overview.tex, deleted once its last topic left;
% shown in a verbatim environment), then the whole of 08_production_considerations.tex below its
% migrate-to line: 11 blocks in order, with only the subsection unnumbered. No other text, table or
% figure was in either file.
% Destinations: def:automated-trading-architecture (blocks 1, 2, 4, 5, 6) and
% meth:execution-decision-support (block 7), owned by execution-trading-costs;
% meth:automated-trading-operational-controls (blocks 8, 9, 10), owned by model-governance-controls;
% def:floor-automated-trading-systems (block 11), owned by execution-market-making;
% block 3 (Broker-Dealer Operations Segmentation) deleted from the compendium and kept only here.
% Sources: Lehalle and Laruelle (alpha-laruelle) secs. 3.1-3.1.3, pp. 193-202; Harris (alpha-harris)
% chap. 7, pp. 147-148, and chap. 27, pp. 543-554; 17 CFR 240.15c3-5 (eCFR, 13 October 2026).
% Unsupported items are listed in the entries' Pitfalls and in notes/migration.md.

% From compendium/legacy/execution-trading-costs/01_execution_and_algorithmic_trading/00_overview.tex (whole file, deleted 2026-10-15).
% Kept in a verbatim environment, exactly as found, so that its \input of the deleted file is shown, not run.
{\scriptsize
\begin{verbatim}
\section{Execution and Algorithmic Trading}


\input{legacy/execution-trading-costs/01_execution_and_algorithmic_trading/08_production_considerations}
\end{verbatim}
}

% From compendium/legacy/execution-trading-costs/01_execution_and_algorithmic_trading/08_production_considerations.tex (whole file, deleted 2026-10-15 once empty).
\subsection*{Production Considerations}

\begin{principle} \hlt{Trading Architecture}\\
Automated trading needs a production architecture that synchronises market data, order gateways, historical databases, models, monitoring, and TCA. The hard part is not one model; it is keeping the decision chain coherent under fragmented venues and live latency.
\end{principle}

\begin{taxonomy} \hlt{Trading Tool Input Segmentation}
\begin{enumerate}[label=\roman*.]
\setlength{\itemsep}{0pt}
\item \hlt{market data}: trades, quotes, depth, timestamps, and book updates, preferably consolidated across venues;
\item \hlt{venue connectivity}: private order-entry messages and public market-data feeds;
\item \hlt{historical data}: raw events, real-time analytics, order messages, fills, and market context;
\item \hlt{models}: impact, risk, alpha, volume, and context models with calibrated parameters;
\item \hlt{calendar}: holidays, auctions, expiries, index events, and expected halt days.
\end{enumerate}
\end{taxonomy}

\begin{taxonomy} \hlt{Broker-Dealer Operations Segmentation}
\begin{enumerate}[label=\roman*.]
\setlength{\itemsep}{0pt}
\item \hlt{front office}: client contact, order receipt, execution, sales, and research;
\item \hlt{back office}: accounts, confirmations, clearing, settlement, custody, and statements;
\item \hlt{risk and treasury}: margin, credit, stock loan, cash management, and firm inventory;
\item \hlt{compliance}: suitability, best execution, surveillance, and regulatory reporting;
\item \hlt{technology}: order routing, market data, drop copies, books, and audit trails.
\end{enumerate}
\end{taxonomy}

\begin{mechanism} \hlt{Analytics Close to the Feed}\\
Spread, effective spread, volatility, book imbalance, and queue-size measures should be computed as soon as raw data arrive. Fast analytics become state variables for execution, monitoring, and later post-trade explanation.
\end{mechanism}

\begin{caveat} \hlt{Connections Are Two-Way}\\
An execution engine sends orders, cancellations, and modifications, then receives acknowledgements, rejects, fills, and public book updates. In dark venues, the private responses can also act as probes of hidden liquidity.
\end{caveat}

\begin{taxonomy} \hlt{Model Component Segmentation}
\begin{enumerate}[label=\roman*.]
\setlength{\itemsep}{0pt}
\item \hlt{formula}: the functional relation between state variables and decisions;
\item \hlt{estimation}: the procedure that calibrates parameters from historical data;
\item \hlt{accuracy check}: the live or post-trade test that decides whether the model is still usable.
\end{enumerate}
\end{taxonomy}

\begin{process} \hlt{Indicator Publication}
\begin{enumerate}[label=\roman*.]
\setlength{\itemsep}{0pt}
\item models and algorithms publish their state as indicators;
\item traders monitor exceptions, concentration, performance, and context shifts;
\item post-trade systems store the same indicators for diagnosis and TCA.
\end{enumerate}
This makes trading algorithms stateful agents in the production stack rather than black boxes.
\end{process}

\begin{process} \hlt{Audit Trail}
\begin{enumerate}[label=\roman*.]
\setlength{\itemsep}{0pt}
\item record order receipt, route, modification, cancellation, execution, allocation, and timestamp;
\item link private order messages to public market state and venue responses;
\item preserve enough context to test best execution, front-running, allocation, and unauthorised trading.
\end{enumerate}
\end{process}

\begin{taxonomy} \hlt{Automated Trading Reliability Segmentation}
\begin{enumerate}[label=\roman*.]
\setlength{\itemsep}{0pt}
\item \hlt{availability}: redundant servers, gateways, power, and network paths;
\item \hlt{integrity}: authenticated users, validated orders, and tamper-resistant logs;
\item \hlt{recoverability}: backups, replayable feeds, drop copies, and disaster-recovery sites;
\item \hlt{capacity}: headroom for message bursts, opens, closes, halts, and data-feed spikes;
\item \hlt{latency}: bounded delays for market data, decisions, order entry, and acknowledgements.
\end{enumerate}
For a trading system, reliability is a market-quality variable, not only an IT variable.
\end{taxonomy}

\begin{caveat} \hlt{Network Failure Creates Position Risk}\\
If order entry, market data, or drop copy fails, the trader may not know whether an order traded, whether a hedge executed, or whether the market moved through a stop.\\
Controls should prefer safe state reconciliation over blind resubmission.
\end{caveat}

\begin{mechanism} \hlt{Automated System Cost Structure}\\
Electronic trading has high fixed costs for software, data networks, security, and recovery, but low marginal labour cost per order. Floor trading has higher human and physical operating costs, but can still be valuable when negotiation, discretion, and soft information dominate raw matching speed.
\end{mechanism}

\end{document}
